% ===========================================================================
% PART 2 -- the estimators, the causal layer, and what is not identified.
% ===========================================================================
\section{The estimators}
\label{sec:est}

\subsection{Why the query family must be declared before anything is measured}

\begin{definition}[Query-relative equivalence]
Given a family $Q$ of queries, $\Nn \sim_Q \Nn'$ iff
$q(\Nn,x) = q(\Nn',x)$ for all $q \in Q$ and all inputs $x$. \emph{Normative
information} is the equivalence class $[\Nn]_Q$, not the text.
\end{definition}

\begin{proposition}[Without $Q$, ``loss'' is not a defined quantity]
\label{prop:Q}
There exist families $Q$ under which every change is total loss, and families
under which no change is any loss. Hence the answer to ``how much was lost'' is
determined entirely by $Q$.
\end{proposition}
\begin{proof}
Take $Q_{\max}$ = all measurable functions of $\Nn$. Then $\Nn\sim_{Q}\Nn'$ iff
$\Nn=\Nn'$ (the identity function is in $Q_{\max}$), so any change destroys the
equivalence entirely. Take $Q_{\min} = \{\text{a constant map}\}$. Then
$\Nn \sim_Q \Nn'$ for all pairs, so no change is a loss. Both families are
admissible in the absence of a stipulation, and they give opposite answers.
\end{proof}

\begin{corollary}
A framework that hides $Q$ inside a supremum has \emph{relocated} the choice, not
made it. Publishing $Q$ before measuring is not a preliminary; it is most of the
result.
\end{corollary}

\subsection{The decision-theoretically correct quantity: normalised regret}

\begin{definition}[Normalised regret]
\label{def:regret}
For a perturbation taking norm $\Nn$ to $\Nn'$, with $a^\star := \argmax_a Y_{\Nn}(a)$
and $a' := \argmax_a Y_{\Nn'}(a)$,
\[
\Reg(\Nn\to\Nn') \;:=\;
\frac{Y_{\Nn}(a^\star) - Y_{\Nn}(a')}{Y_{\Nn}(a^\star) - \min_a Y_{\Nn}(a)}
\;\in\;[0,1].
\]
\end{definition}

\step{Why this and not something else.} $\Reg$ is exactly the Blackwell
deficiency of \cref{thm:garble} restricted to the decision family ``choose one of
these four, scored by the original norm's own utility''. It is not a proxy for
the quantity of interest; on that family it \emph{is} the quantity. The
denominator makes it invariant to the arbitrary scale of $Y$, which
\cref{rem:normalise} showed is not identified.

\begin{definition}[Flip rate]
$\pi := \Prob(a' \neq a^\star)$, estimated by the average of
$\mathbf{1}\{a'\neq a^\star\}$ per \cref{eq:indavg} and \cref{cor:indexp}.
\end{definition}

\begin{lemma}[The flip rate is a coarsening of regret]
\label{lem:coarsen}
$\mathbf{1}\{\text{flip}\} = \mathbf{1}\{\Reg>0\}$, and therefore, exactly,
\[
\E[\Reg] \;=\; \pi \cdot \E[\Reg \mid \text{flip}].
\]
\end{lemma}
\begin{proof}
If $a'=a^\star$ then the numerator of \cref{def:regret} is zero, so $\Reg=0$; if
$a'\neq a^\star$ then $Y_\Nn(a')<Y_\Nn(a^\star)$ by definition of the $\argmax$
(assuming no exact tie), so $\Reg>0$. Hence the indicators coincide. Then
$\E[\Reg] = \E[\Reg \mid \Reg>0]\Prob(\Reg>0) + 0$, by the law of total
expectation, which for a finite space is just splitting the defining sum into the
two events.
\end{proof}

\begin{table}[H]\centering\small
\begin{tabular}{lrrr}
\toprule
operator & $\pi$ & $\E[\Reg]$ & $\E[\Reg\mid\text{flip}]$\\
\midrule
double the target weight & $4.2\%$ & $0.0043$ & $0.1029$\\
delete the target criterion & $5.1\%$ & $0.0067$ & $0.1307$\\
invert the target weight & $11.6\%$ & $0.0294$ & $\mathbf{0.2531}$\\
halve the target weight & $2.3\%$ & $0.0012$ & $0.0532$\\
\bottomrule
\end{tabular}
\caption{$968$ prompts $\times$ $20$ target draws. Target criterion sampled, not
fixed; see \cref{rem:target}.}
\label{tab:regret}
\end{table}

\begin{corollary}[The discarded factor carries information]
By flip rate, invert exceeds delete by $11.6/5.1 = 2.3\times$. By expected
regret, by $0.0294/0.0067 = 4.4\times$. Reporting only the flip rate halves the
measured separation between operators, because inversion not only flips more
often, it flips \emph{worse} -- $25.3\%$ of the available utility range versus
$13.1\%$.
\end{corollary}

\subsection{The flip rate is capped by the margin distribution}

\begin{definition}[Margin]
$m := Y(a_{(1)}) - Y(a_{(2)})$, the gap between the top two responses.
\end{definition}

\begin{proposition}[Margin cap]
\label{prop:cap}
Let $\delta := Y_{\Nn'} - Y_{\Nn}$ be the perturbation of the score vector. Then
$\pi \le \Prob\big(m < 2\|\delta\|_\infty\big)$.
\end{proposition}
\begin{proof}
A flip requires the perturbed score of some $a\neq a^\star$ to exceed that of
$a^\star$, i.e.
$Y(a)+\delta(a) > Y(a^\star)+\delta(a^\star)$, i.e.
$\delta(a)-\delta(a^\star) > Y(a^\star)-Y(a) \ge m$.
Since $|\delta(a)-\delta(a^\star)| \le 2\|\delta\|_\infty$, a flip implies
$m < 2\|\delta\|_\infty$. Taking probabilities and applying monotonicity of
$\Prob$ over the implication gives the bound.
\end{proof}

Measured relative margins $m/(\max_a Y - \min_a Y)$:
\[
p_{05}=0.026,\quad p_{25}=0.144,\quad \text{median}=0.319,\quad p_{75}=0.527,\quad p_{95}=0.800.
\]

\begin{corollary}
An operator whose perturbation is smaller than about a third of the score range
cannot flip half the prompts, \emph{however much normative content it destroys}.
The ceiling on a flip rate is set by the margin distribution, not by the norm.
This is the formal reason \cref{def:regret} is the primary quantity and $\pi$ the
secondary one.
\end{corollary}

\subsection{Target selection is an axis, not a nuisance}
\label{rem:target}

Every criterion-level operator must choose \emph{which} criterion to act on.
Measured effect of that choice on the induced change in $Y$:
\[
\text{highest }|w|: 0.273 \qquad \text{random}: 0.176 \qquad \text{lowest }|w|: 0.041
\]
-- a ratio of $6.6\times$, larger than the differences between operators. Reporting a
single cell would therefore report the selection rule. All numbers in this paper
sweep four selection rules (highest $|w|$, lowest $|w|$, most-rated, random)
$\times$ five seeds, and the seeds are verified to change the draws: only $0.1\%$
of targets are shared across all five.

\subsection{Variance of these estimators}

With $P=968$ prompts and $R=20$ draws each, \cref{thm:deff} applies with measured
$\sigma_b^2 = 0.01361$, $\sigma_w^2 = 0.03677$:
\[
\rho = 0.2701, \qquad \mathrm{DEFF} = 1+19(0.2701) = 6.13, \qquad \sqrt{\mathrm{DEFF}} = 2.48 .
\]

\begin{remark}[Derivation checked against an independent measurement]
A resampling estimate of the same inflation, obtained by recomputing every
pairwise statistic at the cluster level and comparing to the pooled version,
gave $2.79$. Derived $2.48$, measured $2.79$: the residual is that the pairwise
statistic is a correlation rather than a simple mean, for which \cref{thm:deff}
is only approximate.
\end{remark}

\begin{table}[H]\centering\small
\begin{tabular}{lrrrr}
\toprule
operator & $\pi$ & $\se$ & $95\%$ CI & MDE\\
\midrule
double & $4.2\%$ & $0.0036$ & $[3.5,\,4.9]\%$ & $1.00$pp\\
delete & $5.1\%$ & $0.0037$ & $[4.4,\,5.8]\%$ & $1.05$pp\\
invert & $11.6\%$ & $0.0052$ & $[10.6,\,12.6]\%$ & $1.46$pp\\
halve & $2.3\%$ & $0.0026$ & $[1.8,\,2.8]\%$ & $0.72$pp\\
\midrule
\emph{paired} delete $-$ double & $+0.93$pp & $0.0028$ & $z=+3.3$ & $0.79$pp\\
\bottomrule
\end{tabular}
\caption{Standard errors from the $P=968$ cluster means, per \cref{eq:varmean}.
MDE from \cref{thm:mde}.}
\label{tab:power}
\end{table}

\begin{corollary}[Sample size]
By \cref{cor:n} with $\sigma_b = \sqrt{0.01361 + 0.03677/20} = 0.1243$, resolving
a $1$pp difference at $80\%$ power requires
$P = (2.8016 \times 0.1243/0.01)^2 = 1{,}213$ prompts. The release has $968$: the
design is about $20\%$ short of its own headline resolution, and the
delete-versus-double difference clears its MDE only barely.
\end{corollary}

% ===========================================================================
\section{The causal layer: influence is not agreement}
\label{sec:c1}

\subsection{Why a correlational measure cannot answer this}

\begin{definition}[Preservation, causal form]
A distinction reached the output iff \emph{removing it changes the output}.
\end{definition}

A participant whose view happens to match the majority scores perfect alignment
while contributing nothing: remove them and the outcome is unchanged. A
participant who moved the outcome may sit far from it. Alignment and influence
are therefore different functionals, and only the second answers the definition.

\subsection{The Shapley value, derived}

\begin{definition}[Characteristic function]
For prompt $p$ with annotator set $A$ and winner
$w^\star := \argmax_r \sum_{a\in A}\Nn_a(r)$, define for $S \subseteq A$
\[
v(S) \;:=\; \mathbf{1}\Big\{\argmax_r \textstyle\sum_{a\in S}\Nn_a(r) = w^\star\Big\}.
\]
\end{definition}

\begin{definition}[Shapley value]
\[
\phi_i \;:=\; \sum_{S \subseteq A\setminus\{i\}}
\frac{|S|!\,(|A|-|S|-1)!}{|A|!}\;\big[v(S\cup\{i\})-v(S)\big].
\]
\end{definition}

\step{What the weights mean.} Consider all $|A|!$ orderings of the annotators.
The bracket is $i$'s marginal contribution when it arrives after exactly the set
$S$. The number of orderings in which that happens is $|S|!\,(|A|-|S|-1)!$ --
$|S|!$ ways to arrange those before, $(|A|-|S|-1)!$ ways to arrange those after.
Dividing by $|A|!$ makes the weights a probability distribution over orderings.
So $\phi_i$ is \textbf{the average marginal contribution of $i$ over a uniformly
random arrival order}, which is exactly the causal definition above, averaged
over every context in which $i$ could arrive.

\step{Estimation.} The sum has $2^{|A|-1}$ terms. We sample $200$ random
orderings per prompt and average the realised marginal contributions, which is an
unbiased estimator of $\phi_i$ by construction (each ordering is drawn from the
same uniform distribution the weights encode).

\subsection{Result: the pre-registered kill does not fire}

\begin{table}[H]\centering\small
\begin{tabular}{lr}
\toprule
$1{,}865$ (annotator, prompt) pairs & \\
\midrule
$\Corr(\text{Shapley influence},\ \text{cosine alignment})$ & $+0.3040$\\
bootstrap $95\%$ CI (\cref{prot:boot}) & $[+0.2566,\,+0.3520]$\\
pre-registered kill: correlation $\approx 1.0$ & \textbf{does not fire}\\
\bottomrule
\end{tabular}
\end{table}

\begin{table}[H]\centering\small
\begin{tabular}{lrr}
\toprule
alignment quintile & mean alignment & mean influence\\
\midrule
Q1 & $-0.0416$ & $0.0083$\\
Q2 & $+0.8164$ & $0.1055$\\
Q3 & $+0.9489$ & $\mathbf{0.1325}$\\
Q4 & $+0.9861$ & $0.1262$\\
Q5 & $+0.9973$ & $0.1262$\\
\bottomrule
\end{tabular}
\caption{Influence rises steeply from Q1 to Q3 and then \emph{flattens}.}
\label{tab:quintile}
\end{table}

\mesotakeaway{Beyond moderate agreement, agreeing more buys no influence at all.
And the off-diagonal cells are populated: $0.75\%$ of pairs are top-quintile
aligned with bottom-quintile influence, $3.11\%$ are bottom-aligned with
top-quintile influence. ``Agreement is not preservation'' is here a count of real
people, not a slogan.}

\subsection{An independent confirmation from a different operator}

Dropping the \emph{most conforming} annotator changes the decision on
$5.8\%\,[4.3,7.1]$ of prompts; dropping the \emph{largest dissenter}, on
$2.1\%\,[1.1,3.2]$. The intervals do not overlap and the sign is stable across
all three judges. The dissenter was already not influencing the outcome -- the
same statement as the flat top of \cref{tab:quintile}, reached by a different
route.

% ===========================================================================
\section{What is not identified, with proofs}
\label{sec:notid}

\begin{proposition}[Contamination and redundancy are not separable here]
\label{prop:contam}
Let $\alpha_p$ be the alignment between a prompt's weighted-criteria score and
the crowd's ranking, and $s_p$ the sensitivity of its decision to perturbing one
criterion. Two models,
\begin{itemize}
\item[$C$:] criteria written \emph{after} seeing responses are individually more
      decisive, so a high-$\alpha$ rubric absorbs single-criterion perturbations;
\item[$R$:] a high-$\alpha$ rubric is one whose criteria agree with each other,
      and a redundant rubric is insensitive to any single member,
\end{itemize}
imply the same sign for $\Corr(\alpha,s)$, and the observed data distribution is
a function of $(\alpha,s)$ alone. The likelihoods coincide; the models are not
identified.
\end{proposition}

\step{What is identified: a bound.} Splitting prompts at the median of $\alpha$:

\begin{table}[H]\centering\small
\begin{tabular}{lrrr}
\toprule
operator & low $\alpha$ & high $\alpha$ & ratio\\
\midrule
double & $5.1\%$ & $2.8\%$ & $0.55\times$\\
delete & $5.6\%$ & $4.6\%$ & $0.82\times$\\
invert & $13.1\%$ & $9.9\%$ & $0.75\times$\\
\bottomrule
\end{tabular}
\end{table}

\begin{corollary}
Every flip rate reported here varies by up to $1.8\times$ across the range of
$\alpha$. That bound belongs on every such number. Separating $C$ from $R$
requires response-blind elicitation and cannot be done at this site.
\end{corollary}

\begin{proposition}[The probe bound is an identification limit]
By \cref{thm:probe}, at most two parameters per person per prompt are identified
from the ratings. Everything else in the criterion text is unidentified from the
released quantitative data.
\end{proposition}

\begin{proposition}[Compilation's three actions are confounded]
Selection, rewriting and weight-discarding occur in one step with no lineage.
Their individual contributions to the $29.4\%$ disagreement are not identified.
\end{proposition}

\begin{proposition}[$Y$ does not exist]
The chain in \cref{eq:chain} terminates at $D$. $Y$ -- what a system trained under
this standard does -- is not in the release, and no arithmetic recovers it. Any
end-to-end number quoted to $Y$ is fabricated.
\end{proposition}

% ===========================================================================
\section{Findings}
\label{sec:findings}

Each finding ends with its confidence card. A missing entry reads
\textsc{uncomputed}; none is omitted.

\subsection{F1 --- the elicitation captured consensus, not the individual}

$\cos(\Nn_i,\text{own ranking}) = 0.3928$ versus
$\cos(\Nn_i,\text{stranger's ranking}) = 0.3097$; difference $+0.0831$.

\begin{table}[H]\centering\footnotesize
\begin{tabular}{ll}
\toprule
$n_{\text{eff}}$ & $968$ prompt clusters ($14{,}764$ person--prompt rows)\\
MDE & $0.0187$ in cosine, from \cref{thm:mde} at the observed $\se$\\
effect / floor & $0.0831 / 0.0067 = 12.4$\\
CI width / $|$eff$|$ & $0.0262/0.0831 = 0.32$\\
spec\_survival & \textsc{uncomputed} (single specification)\\
seed spread / $|$eff$|$ & not applicable (no stochastic step)\\
held out & \textsc{absent}\\
instrument & routes through the judge; ceiling $0.720$ (§\ref{sec:instrument})\\
prior art & none located; the release documents no such comparison\\
multiplicity & $1$ test in this family\\
\bottomrule
\end{tabular}
\end{table}

\subsection{F2 --- no finite weight encodes a veto}

\cref{thm:veto} and \cref{thm:vetoinf}; required $M$ up to $26.54$ against a scale
maximum contribution of $10$.

\begin{table}[H]\centering\footnotesize
\begin{tabular}{ll}
\toprule
$n_{\text{eff}}$ & $161$ prompts with a partial veto; $42$ binding\\
MDE & not applicable --- this is a \textbf{derivation}, not a test\\
effect / floor & not applicable\\
CI width / $|$eff$|$ & not applicable\\
spec\_survival & the theorem holds for every score construction\\
seed spread & not applicable (deterministic)\\
held out & the \emph{maximum} is what carries it; a held-out set can only raise it\\
instrument & $\Delta$ is computed from the judge tensor; the \emph{theorem} is not\\
prior art & the incompatibility of vetoes with scoring rules is classical; the\\
 & \emph{measured} $\Delta$ distribution on this release is not\\
multiplicity & none --- a single derived statement\\
\bottomrule
\end{tabular}
\end{table}

\begin{remark}[Labelled as a derivation, per the arithmetic trap]
\cref{thm:vetoinf} could not have come out otherwise: it follows from
\cref{thm:veto} plus the unboundedness of $\Delta$. It is stated as a theorem,
its assumption is named ($\sat$ can approach its bounds independently across
responses), and it is not counted as evidence in the sense a measurement is.
\end{remark}

\subsection{F3 --- comparative measurements are instrument-invariant; absolute ones are not}

\cref{tab:gaugea}, \cref{tab:gaugeb}.

\begin{table}[H]\centering\footnotesize
\begin{tabular}{ll}
\toprule
$n_{\text{eff}}$ & $968$ prompts (three instruments); $300$ for the reordering\\
MDE & $\rho$ resolution $\approx 0.15$ at $11$ operators\\
effect / floor & orderings $0.963$--$1.000$ against a chance $\rho$ of $0$\\
CI width / $|$eff$|$ & bootstrap intervals in \cref{tab:power}\\
spec\_survival & $3$ of $3$ instruments preserve the ordering\\
seed spread / $|$eff$|$ & $<0.1$ across $5$ seeds of the target draw\\
held out & the reordered run is a genuine held-out instrument\\
instrument & this finding \emph{is} about the instrument, by construction\\
prior art & judge-sensitivity of LLM evaluation is known; the\\
 & comparative/absolute split measured on this object is not\\
multiplicity & $3$ instruments $\times$ $11$ operators; the reported statistic\\
 & is the ordering, one per instrument\\
\bottomrule
\end{tabular}
\end{table}

\subsection{F4 --- influence and alignment are distinct (C1 survives)}

$\Corr = +0.3040\,[+0.2566,+0.3520]$, with the quintile profile of
\cref{tab:quintile}.

\begin{table}[H]\centering\footnotesize
\begin{tabular}{ll}
\toprule
$n_{\text{eff}}$ & $968$ prompt clusters ($1{,}865$ annotator--prompt rows)\\
MDE & $0.065$ in correlation, at the bootstrap $\se$\\
effect / floor & $0.304 / 0.024 = 12.7$\\
CI width / $|$eff$|$ & $0.095/0.304 = 0.31$\\
spec\_survival & \textsc{uncomputed}; one characteristic function tested\\
seed spread / $|$eff$|$ & $<0.02$ over the $200$-permutation Shapley estimator\\
held out & \textsc{absent}\\
instrument & routes through the judge; ceiling $0.720$\\
prior art & Shapley-versus-alignment has not been reported on this object\\
multiplicity & $1$ pre-registered test\\
\bottomrule
\end{tabular}
\end{table}

\subsection{F5 --- one bit per criterion reproduces $89.8\%$ of decisions}

\cref{tab:rd}, cross-checked against \cref{tab:cardinal}.

\begin{table}[H]\centering\footnotesize
\begin{tabular}{ll}
\toprule
$n_{\text{eff}}$ & $968$ prompts\\
MDE & $\pm 1.0$pp at the observed $\se$\\
effect / floor & the quantisation is deterministic; no resampling floor applies\\
CI width / $|$eff$|$ & \textsc{uncomputed}\\
spec\_survival & agrees with the independent sign-only test to $1.5$pp\\
seed spread & not applicable (deterministic)\\
held out & \textsc{absent}\\
instrument & judge-dependent in level; the \emph{bit curve} is not tested across judges\\
prior art & none located for this object\\
multiplicity & $5$ quantisation levels, reported in full\\
\bottomrule
\end{tabular}
\end{table}

% ===========================================================================
\section{What did not survive}
\label{sec:exclusion}

A paper with no retractions is a paper whose ontology has no word for error. All
entries are cross-referenced to \texttt{supplement/exclusion\_log.csv}.

\begin{table}[H]\centering\footnotesize
\begin{tabularx}{\textwidth}{P{4.2cm}P{4.6cm}Y}
\toprule
withdrawn claim & what killed it & status\\
\midrule
``the operator design has rank $2$, so all remaining dimensions require the
judge'' & the rank statistic sits inside the span of three defensible nulls
($12.1$, $17.0$, $18.0$) with the observation at $13$; the null
\emph{construction} moved it more than the data & rank retired as a statistic\\
``strengthening a criterion moves the decision, removing one does not'' & direct
measurement: deleting flips $8.8\%$, doubling $4.2\%$. The asymmetry was an
artefact of normalising the score \emph{after} mutating and then differencing two
separately normalised vectors & withdrawn, sign and all\\
``minority-sensitive aggregation changes half the outcomes'' & the sham control
(same people, matched moments, response pattern destroyed) disagrees $61.4\%$
versus the real $49.4\%$ & withdrawn; the residual runs the other way\\
``aggregation is a $10\times$ larger lever than the criterion layer'' & the
aggregators were indexed over criteria rather than people & magnitude withdrawn;
ordering survives at $4.3\times$\\
``the flip rate is confounded by decision margin'' & measured $Q_1/Q_4$ ratios of
$1.0$, $1.0$, $1.1$ & defect opened by us and closed by its own measurement\\
``wide margins arise because agreeing criteria carry larger weights'' &
$\Corr(\text{margin},\overline{|w|}) = +0.002\,[-0.067,+0.059]$ & mechanism
refuted; the flatness stands unexplained\\
\bottomrule
\end{tabularx}
\end{table}

\begin{remark}[The failure mode behind most of these]
Six of the six are the same shape: a verdict decided by a constant the author
chose -- one null instead of a range, one reference instead of a class, a
two-branch template for a three-outcome test, a $0.15$ cutoff. The repair adopted
here is structural: \emph{every verdict is a printed comparison against the range
of defensible references, never a boolean computed from a threshold}.
\end{remark}

% ===========================================================================
\section{Falsification and contradiction handling}
\label{sec:falsify}

Stated before the claims are defended.

\begin{enumerate}
\item \textbf{F1 dies} if a response-blind elicitation shows
      $\cos(\Nn_i,\text{own})$ far exceeding $\cos(\Nn_i,\text{stranger})$. The
      present gap would then be an artefact of post-hoc writing rather than a
      property of elicitation.
\item \textbf{F2 dies} if a bounded $M$ is exhibited that reproduces
      constraint behaviour across candidate sets. \cref{thm:vetoinf} forbids it
      for unbounded $\sat$; a bounded-satisfaction variant with a stated bound
      would restrict the theorem's scope.
\item \textbf{F3 dies} if a fourth instrument reverses the operator ordering
      ($\rho<0.8$). Three have not.
\item \textbf{F4 dies} if $\Corr(\text{influence},\text{alignment})\to 1$ under a
      different characteristic function. Only one has been tested; this is the
      weakest of the five.
\item \textbf{F5 dies} if the bit curve differs materially under another judge.
      Not tested.
\item \textbf{The framing itself dies} if a person's induced ordering is unstable
      across response sets. Then what was elicited is a property of the probe,
      not of the person, and nothing downstream measures values at all. This is
      the deepest available test and it requires two candidate sets per person,
      which this release does not have. \cref{cor:iia} ($13.4\%$ IIA violations)
      is a lower bound on the concern, not a resolution of it.
\end{enumerate}

% ===========================================================================
\section{Discussion}

\subsection{The claim, and its novelty boundary}

\textbf{What is not new.} Blackwell's theorem, Arrow's impossibility, Sen's
liberal paradox, the Shapley value, Spearman--Brown, the design effect,
attenuation from measurement error. The incompatibility between vetoes and
scoring rules is classical social choice. Judge-sensitivity of language-model
evaluation is documented. None of these is claimed here.

\textbf{What is new, and only this.} (i) The measured $\Delta$ distribution that
turns the veto incompatibility into a \emph{quantitative} statement about a
specific release, with the counterexample present in the data
(\cref{tab:vetoM}). (ii) The self-versus-stranger decomposition showing that the
individual signal in this elicitation is $0.083$ of cosine (\cref{tab:selfother}).
(iii) The comparative/absolute split under the gauge test (§\ref{sec:gauge}).
(iv) The bit-rate of the weight scale (\cref{tab:rd}). (v) The demonstration that
the weighted-criteria score agrees with the Condorcet winner only $68.4\%$ of the
time while every ranking rule agrees $97$--$98\%$.

\subsection{Why method relativity is not relativism}

Every quantity here is relative to a declared query family
(\cref{prop:Q}), a declared decision family (\cref{def:regret}), and a named
instrument (§\ref{sec:instrument}). That is not an admission that the numbers are
arbitrary; it is the statement that they are \emph{typed}. A claim inherits the
scope of the family it was computed against and may not be quoted outside it.
The gauge test (§\ref{sec:gauge}) is what makes the typing empirical rather than
rhetorical: it shows exactly which claims move when the instrument moves.

% ===========================================================================
\section{Limitations and epistemic status}
\label{sec:limits}

``We did not run $X$'' and ``$X$ cannot be run at this site'' are different
sentences and are not merged.

\subsection{Structural --- cannot be run here}

\begin{enumerate}
\item \textbf{Response-blind source norms.} Every criterion was written after
      seeing the responses. \cref{prop:contam} shows contamination and redundancy
      are not separable without them.
\item \textbf{A second candidate set per person.} Required to test
      \cref{thm:vetoinf} out of sample and to resolve the probe-dependence
      question of §\ref{sec:falsify}(6).
\item \textbf{Typed fields.} $82.3\%$ of criteria carry none
      (\cref{tab:payload}); scope and exception loss are unstudiable here.
\item \textbf{Compilation lineage.} Absent, so §\ref{sec:compile}'s three actions
      stay confounded.
\item \textbf{$Y$.} Not in the release.
\item \textbf{Cross-dataset / cross-domain.} One release exists.
\end{enumerate}

\subsection{Practical --- could be run, was not}

\begin{enumerate}
\item Held-out partitions for F1, F4, F5.
\item A specification curve over characteristic functions for the Shapley
      estimator; one was tested.
\item The bit curve of \cref{tab:rd} across judges.
\item Two-way clustering (\cref{thm:cgm}) on the person-level quantities; only
      prompt-level clustering was applied, and annotators repeat across prompts.
\end{enumerate}

\subsection{Quantitative bounds that apply to every number here}

\begin{table}[H]\centering\small
\begin{tabular}{ll}
\toprule
bound & value\\
\midrule
judge attenuation ceiling & $0.741$\\
panel reliability ceiling & $0.9722$\\
compounded ceiling on any correlation & $0.720$\\
contamination-proxy variation in every flip rate & up to $1.8\times$\\
absolute winner claims: instrument share & $\approx 40\%$\\
absolute winner claims: presentation-order share & $8.0$pp\\
rubric corpus never joined & $1.8\%$\\
design shortfall for $1$pp resolution & $968$ of $1{,}213$ prompts\\
pairwise-dependence screen: MDE at full $n$ & $|\phi| = 0.153$ against median $0.124$\\
\bottomrule
\end{tabular}
\end{table}

\begin{remark}[The last row read honestly]
The median pairwise association in the operator grid is \emph{below} the MDE at
the largest available sample. The pairwise table is therefore a screen for
\emph{strong} dependence and is worded that way throughout; its nulls are
underpowered rather than informative.
\end{remark}

% ===========================================================================
\section{What a design that can measure this requires}
\label{sec:design}

Each row is derived from a specific result above, not proposed.

\begin{table}[H]\centering\footnotesize
\begin{tabularx}{\textwidth}{P{4.6cm}P{3.4cm}Y}
\toprule
requirement & derived from & unlocks\\
\midrule
response-blind source norms & \cref{prop:contam} & the only route separating
contamination from redundancy\\
$\ge 2$ candidate sets per person & \cref{thm:vetoinf}, \cref{cor:iia} & tests
whether $M$ transfers; quantifies candidate-set overfitting\\
typed fields: veto, scope, exception, priority as separate columns &
§\ref{sec:type}, \cref{tab:payload} & the index mismatch is a missing column, not
a hard problem\\
interval-scale validation, or paired comparisons instead & \cref{tab:cardinal} &
$11.7\%$ of decisions currently rest on an unestablished assumption\\
$k>4$ responses & \cref{thm:probe} & raises the probe bound from $2$ to $k-2$\\
$\ge 1{,}213$ prompts & \cref{cor:n} & $1$pp resolution\\
substantially more than $968$ for pairwise structure & §\ref{sec:limits} &
the typical association is currently undetectable\\
$\ge 3$ executors & §\ref{sec:gauge} & already available; the gauge is established\\
compilation receipt (lineage) & §\ref{sec:compile} & separates selection,
rewriting, weight-discarding\\
declared $Q$ and declared $\Agg$ & \cref{prop:Q}, \cref{rem:threelosses} &
without them ``loss'' is undefined and ``defect'' is inseparable from ``social
choice''\\
\bottomrule
\end{tabularx}
\end{table}

% ===========================================================================
\section{Conclusion}

The chain of \cref{eq:chain} loses something at every arrow, and the losses are
not of one kind.

\begin{itemize}
\item \textbf{$\succeq_i \to \Nn_i$.} An unbounded relation is projected onto two
      free parameters (\cref{thm:probe}) -- forced, not recoverable. Within that,
      a person's own criteria recover $39\%$ of the direction of their own choice,
      and the individual component of that is $0.083$ of cosine. The type
      structure is empty: $82.3\%$ of criteria carry no condition, exception,
      priority, hedge or prohibition. And $11.7\%$ of decisions depend on an
      interval reading of a scale that was never validated as one.
\item \textbf{$\Nn_i \to G$.} Aggregation costs $0.1298$ of normalised utility
      (\cref{thm:garble}) -- a social choice, but one made silently because
      $\Agg$ was never declared. The weighted-criteria route agrees with the
      Condorcet winner only $68.4\%$ of the time while ranking rules agree
      $97$--$98\%$. $13.4\%$ of pairwise relations flip when the option set
      changes. And $26.1\%$ of the time the utilitarian winner is a response
      somebody called unacceptable.
\item \textbf{Type.} No finite weight encodes a veto (\cref{thm:vetoinf}). This
      is the strongest statement in the paper and it required no instrument.
\item \textbf{$G \to C$.} Compilation changes $29.4\%$ of decisions, with its
      three actions confounded. But the discarded magnitudes were worth about one
      bit: sign alone reproduces $89.8\%$.
\item \textbf{$C \to D$.} The executor's reliability is $0.5497$, capping every
      correlation at $0.741$. Changing the reader changes $40\%$ of winners;
      reordering the page changes $8\%$. Comparative claims survive all of it;
      absolute claims do not.
\item \textbf{$D \to Y$.} Does not exist in the release.
\end{itemize}

\mesotakeaway{The losses inside the elicitation are larger than every loss the
pipeline adds. The most defensible statement this object supports is a
\emph{theorem} about what a scalar scale cannot encode, and the most consequential
one is that the next dataset should buy \textbf{fields and response-blindness},
not finer weights or more prompts.}
